% =====================================================================
%  NATURAL SCIENCES TRIPOS  -  PART II PHYSICS
%  ADVANCED AND QUANTUM OPTICS  -  Examination Paper
%
%  Compile with pdflatex.  No external packages beyond the standard
%  TeX Live distribution are required.
% =====================================================================
\documentclass[11pt]{article}

\usepackage[a4paper,top=2.4cm,bottom=2.4cm,left=2.6cm,right=2.6cm]{geometry}
\usepackage{amsmath,amssymb}
\usepackage{graphicx}
\usepackage{tikz}
\usepackage{enumitem}
\usepackage{fancyhdr}
\usepackage{mathptmx}   % Times-like font, in keeping with Tripos papers

% ---------------------------------------------------------------------
%  Layout helpers
% ---------------------------------------------------------------------
\setlength{\parindent}{0pt}
\setlength{\parskip}{6pt}

% Marks in the right margin -------------------------------------------
\newcommand{\qmarks}[1]{\hfill\makebox[0pt][r]{\hspace{1.5em}\textbf{[#1]}}}

% Question numbering --------------------------------------------------
\newcounter{question}
\newcommand{\question}{\par\vspace{4pt}\refstepcounter{question}%
  \noindent\makebox[1.6em][l]{\textbf{\thequestion}}}

% Part labelling (a), (b), ... ----------------------------------------
\newlist{parts}{enumerate}{1}
\setlist[parts]{label=(\alph*),leftmargin=2.2em,labelsep=0.6em,
                itemsep=4pt,topsep=4pt}

% Page style ----------------------------------------------------------
\pagestyle{fancy}
\fancyhf{}
\rfoot{\textit{\small (TURN OVER}}
\cfoot{}
\lfoot{}
\renewcommand{\headrulewidth}{0pt}
\renewcommand{\footrulewidth}{0pt}

% Section banner ------------------------------------------------------
\newcommand{\sectionbanner}[2]{%
  \vspace{6pt}\noindent{\large\textbf{#1}}\\[-2pt]
  \noindent\textit{#2}\par\vspace{4pt}}

\begin{document}

% =====================================================================
%  COVER SHEET
% =====================================================================
\thispagestyle{empty}

\vspace*{0.4cm}
{\large\textbf{NATURAL SCIENCES TRIPOS} \hfill Part II}

\vspace{0.3cm}
\hrule
\vspace{0.4cm}

\noindent Wednesday 10th June 2026 \hfill 09:00 to 11:00

\vspace{0.3cm}
\hrule
\vspace{0.8cm}

{\large\textbf{PHYSICS (5)}}\\[2pt]
{\large\textbf{PHYSICAL SCIENCES: HALF SUBJECT PHYSICS (5)}}

\vspace{0.5cm}
{\large\textbf{ADVANCED AND QUANTUM OPTICS}}

\vspace{0.8cm}
\textit{Candidates offering this paper should attempt \textbf{all six} questions
(four from Section A and two from Section B).}

\vspace{0.5cm}
\textit{The approximate number of marks allocated to each question or part of a
question is indicated in the right margin. This paper contains 4 sides,
including this coversheet, and is accompanied by a handbook giving values of
constants and containing mathematical formulae which you may quote without
proof.}

\vspace{1.0cm}
\begin{tabular}{@{}p{6.4cm}p{6.4cm}@{}}
\textbf{STATIONERY REQUIREMENTS} & \textbf{SPECIAL REQUIREMENTS}\\[4pt]
$2\times8$-page answer booklet      & Mathematical Formulae Handbook\\
Rough workpad                       & Approved calculator allowed\\
Yellow master coversheet            & \\
\end{tabular}

\vspace{1.4cm}
\begin{center}
\fbox{\fbox{\parbox{9.2cm}{\centering
You may not start to read the questions printed on the subsequent pages of
this question paper until instructed that you may do so by the Invigilator.}}}
\end{center}

\vfill
\noindent\copyright\ 2026 University of Cambridge

\newpage

% =====================================================================
%  SECTION A
% =====================================================================
\setcounter{page}{2}
\begin{center}\textbf{2}\end{center}

\sectionbanner{SECTION A}{Attempt \textbf{all four} questions in this Section.
Answers should be concise and relevant formulae may be assumed without proof.}

% ----------------------------- Q1 ------------------------------------
\question
A collimated Gaussian beam of wavelength $\lambda = 1.0\ \mu$m and waist radius
$W_0 = 0.50$~mm propagates in free space. State what is meant by the
\textit{Rayleigh range} of the beam, evaluate it for this beam, and hence
estimate the beam diameter at a distance of $20$~m from the waist. Comment
briefly on why a larger waist gives a more collimated beam.
\qmarks{4}

% ----------------------------- Q2 ------------------------------------
\question
A second-order nonlinear crystal is illuminated by a single monochromatic field
of angular frequency $\omega$. Explain physically why the nonlinear
polarisation contains components at both $2\omega$ and at zero frequency, and
name the two corresponding optical processes. State the condition on the
crystal symmetry that is required for these processes to occur, and explain
why it is not satisfied in a centrosymmetric medium.
\qmarks{4}

% ----------------------------- Q3 ------------------------------------
\question
A symmetric two-mirror optical resonator is constructed from two identical
concave mirrors of radius of curvature $R$, separated by a distance $d$.
Derive the stability (confinement) condition for this resonator in terms of
$d$ and $R$, and identify the values of $d/|R|$ that correspond to the planar,
confocal and concentric configurations. State which of these is the most
robust against mirror misalignment, and why.
\qmarks{4}

% ----------------------------- Q4 ------------------------------------
\question
A single two-level atom is driven on resonance by a classical light field with
resonant Rabi frequency $\Omega_0$. Sketch the population of the excited state
as a function of time (i) in the absence of spontaneous emission, and (ii) when
spontaneous emission at rate $\gamma$ is included, with $\gamma \approx
0.2\,\Omega_0$. In each case state the behaviour at long times and explain
physically the difference between the two cases.
\qmarks{4}

\newpage
\begin{center}\textbf{3}\end{center}

% =====================================================================
%  SECTION B
% =====================================================================
\sectionbanner{SECTION B}{Attempt \textbf{both} questions from this Section.}

% ----------------------------- Q5 ------------------------------------
\question
This question concerns optical coherence and its use in interferometric
imaging.

A quasi-monochromatic light source of mean frequency $\nu_0$ emits a partially
coherent plane wave whose complex degree of temporal coherence may be written
$g(\tau) = g_a(\tau)\exp(j2\pi\nu_0\tau)$, where $g_a(\tau)$ is a slowly
varying envelope.

\begin{parts}

\item Define the \textit{coherence time} $\tau_c$ and the \textit{coherence
length} $\ell_c$ of the source, and state how $\tau_c$ is related to the
spectral width $\Delta\nu_c$ of the light.
\qmarks{3}

\item The wave is sent into a Michelson interferometer; one arm is lengthened
relative to the other so as to introduce a time delay $\tau$ between the
recombining fields, which have equal intensities $I_0$. Show that the detected
intensity may be written
\[
  I(\tau) = 2I_0\bigl[\,1 + |g(\tau)|\cos\varphi(\tau)\,\bigr],
\]
identifying $\varphi(\tau)$. Sketch and annotate $I(\tau)$ for (i) a perfectly
monochromatic wave and (ii) a quasi-monochromatic wave, and explain how the
interferogram allows $|g(\tau)|$ to be measured.
\qmarks{5}

\item A light-emitting diode has a centre wavelength of $\lambda_0 = 800$~nm
and a spectral width of $\Delta\lambda = 25$~nm. Estimate its coherence length.
Explain why a short coherence length is desirable for an interferometric
imaging technique that profiles a multilayered sample.
\qmarks{4}

\item Optical coherence tomography (OCT) uses such a source to image a sample
consisting of several weakly reflecting boundaries at depths giving delays
$\tau_i$. Explain qualitatively why each boundary produces a distinct burst of
interference fringes in the interferogram, and state how the axial separation
of two layers is extracted from the measurement.
\qmarks{4}

\item State one reason why the axial resolution of OCT is, unlike that of a
conventional microscope, independent of the numerical aperture of the imaging
optics.
\qmarks{4}

\end{parts}

\newpage
\begin{center}\textbf{4}\end{center}

% ----------------------------- Q6 ------------------------------------
\question
This question concerns the quantised single-mode electromagnetic field and the
states of light it supports.

A single mode of the electromagnetic field is described by the Hamiltonian
$\hat{H} = \hbar\omega(\hat{a}^{\dagger}\hat{a} + \tfrac12)$, where $\hat{a}$
and $\hat{a}^{\dagger}$ obey $[\hat{a},\hat{a}^{\dagger}] = 1$. The electric
field operator may be taken as
$\hat{E}_x \propto (\hat{a} + \hat{a}^{\dagger})$.

\begin{parts}

\item The number (Fock) states $|n\rangle$ are the eigenstates of
$\hat{n} = \hat{a}^{\dagger}\hat{a}$. Using
$\hat{a}|n\rangle = \sqrt{n}\,|n-1\rangle$ and
$\hat{a}^{\dagger}|n\rangle = \sqrt{n+1}\,|n+1\rangle$, show that
$\langle n|\hat{E}_x|n\rangle = 0$ for all $n$, while
$\langle n|\hat{E}_x^{2}|n\rangle \propto (n + \tfrac12)$. Comment on why a
Fock state, even with large $n$, does not behave like a classical light field.
\qmarks{5}

\item A \textit{coherent state} $|\alpha\rangle$ is defined as the eigenstate of
the annihilation operator, $\hat{a}|\alpha\rangle = \alpha|\alpha\rangle$ with
$\alpha\in\mathbb{C}$. Given that
\[
  |\alpha\rangle = e^{-|\alpha|^{2}/2}\sum_{n=0}^{\infty}
  \frac{\alpha^{n}}{\sqrt{n!}}\,|n\rangle,
\]
show that the probability of detecting $n$ photons is Poissonian, and hence
that the mean photon number is $\bar{n} = |\alpha|^{2}$ and the variance is
$(\Delta n)^{2} = \bar{n}$.
\qmarks{5}

\item Explain why a coherent state, unlike a Fock state, has a non-zero
expectation value of the field, and state how the fractional uncertainty in the
photon number, $\Delta n/\bar{n}$, behaves as $|\alpha| \to \infty$. Why does
this justify regarding the coherent state as the quantum description of an
ideal laser beam?
\qmarks{4}

\item A single photon, in the state $|0\rangle_{0}\,|1\rangle_{1}$, is incident
on one input port of a lossless 50:50 beam splitter, the other input port being
in the vacuum state. Using the beam-splitter relation
$\hat{a}^{\dagger}_{1} = \tfrac{1}{\sqrt2}
(i\,\hat{a}^{\dagger}_{2} + \hat{a}^{\dagger}_{3})$,
find the output state. Comment on the role played by the vacuum entering the
unused port.
\qmarks{3}

\item Two indistinguishable single photons are now incident, one on each input
port of the same beam splitter. State, with a brief justification, what is
observed in a coincidence measurement at the two outputs, and explain how this
effect is used to test whether two photons are truly indistinguishable.
\qmarks{3}

\end{parts}

\vspace{1.2cm}
\begin{center}\textbf{END OF PAPER}\end{center}

\thispagestyle{empty}
\vfill
\noindent\copyright\ 2026 University of Cambridge

\end{document}
